\subsection{Binary witness tables and lossless field packing}\label{sec:packing:tables}
Fix integers $\ell\geq1$ and $h\geq0$, and put
\[
 n:=2^\ell,\qquad L:=\F_{2^n},\qquad q:=|L|=2^n,\qquad
 I:=I_\ell=\{0,1\}^\ell,\qquad J:=I_h=\{0,1\}^h.
\]
Choose $\alpha\in L$ of degree $n$ over $\F_2$ and use the
kernel evaluation basis $\beta_v=K_v^{(\ell)}(\alpha)$,
$v\in I$, and trace-dual basis $\eta_v$ from
Section~\ref{sec:binary:coordinates}.
The equality weights are the polynomials
$\operatorname{eq}_d$ from Chapter~1, interpreted as
the constant $1$ when $d=0$.
The packed-ring-switching construction and its history
are described in~\cite{DiamondPosen2024}.
We derive each required identity rather than importing
a protocol-level correctness statement.

\begin{definition}[Boolean witness and its field embedding]\label{def:pack:witness}
A \emph{Boolean witness} is a function
$w:I\times J\to\F_2$. We identify the prime subfield
$\F_2\subseteq L$ with $\{0,1\}\subseteq L$
when forming products or sums in $L$.
For $v\in I$ and $j\in J$ write $w_{v,j}:=w(v,j)$.
\end{definition}

\begin{definition}[Packing a Boolean column]\label{def:pack:pack}
For $a=(a_v)_{v\in I}\in\F_2^I$ define
\[
 \operatorname{Pack}(a):=\sum_{v\in I}\beta_v a_v\in L.
\]
For $x\in L$ define
$\operatorname{Unpack}(x):=(d_v(x))_{v\in I}\in\F_2^I$,
where $d_v$ are the unique coordinates in the basis
$(\beta_v)$.
\end{definition}

\begin{lemma}[Unpack after Pack]\label{lem:pack:unpack-pack}
For every $a\in\F_2^I$,
$\operatorname{Unpack}(\operatorname{Pack}(a))=a$.
\end{lemma}
\begin{proof}
The element $\operatorname{Pack}(a)$ is by definition
$\sum_v\beta_va_v$. By uniqueness of coordinates
in the basis $(\beta_v)$,
$d_u(\operatorname{Pack}(a))=a_u$ for every $u\in I$.
These equalities are the equality of the two vectors.
\end{proof}

\begin{lemma}[Pack after Unpack]\label{lem:pack:pack-unpack}
For every $x\in L$,
$\operatorname{Pack}(\operatorname{Unpack}(x))=x$.
\end{lemma}
\begin{proof}
The defining coordinate expansion is
$x=\sum_v\beta_vd_v(x)$. Its right side is
$\operatorname{Pack}(\operatorname{Unpack}(x))$.
\end{proof}

\begin{lemma}[Packing is binary-linear]\label{lem:pack:linear}
For $a,b\in\F_2^I$ and $e\in\F_2$,
$\operatorname{Pack}(a+eb)=
\operatorname{Pack}(a)+e\operatorname{Pack}(b)$.
The map $\operatorname{Unpack}$ is also $\F_2$-linear.
\end{lemma}
\begin{proof}
Distribute the finite sum
$\sum_v\beta_v(a_v+eb_v)$, using
$e\in\F_2$. Each coordinate map $d_v$ is
$\F_2$-linear by Lemma~\ref{lem:binary:extract},
so the coordinatewise map $\operatorname{Unpack}$ is.
\end{proof}

\begin{theorem}[Lossless Boolean packing]\label{thm:pack:bijection}
$\operatorname{Pack}:\F_2^I\to L$ and
$\operatorname{Unpack}:L\to\F_2^I$ are mutually inverse
$\F_2$-linear isomorphisms.
\end{theorem}
\begin{proof}
Lemmas~\ref{lem:pack:unpack-pack} and
\ref{lem:pack:pack-unpack} show that the maps are
mutually inverse. Lemma~\ref{lem:pack:linear}
proves their $\F_2$-linearity.
\end{proof}

\begin{definition}[Packed witness and its column coordinates]\label{def:pack:columns}
For each $j\in J$ define the packed field element
\[
 \widehat w_j:=\operatorname{Pack}
                    \bigl((w_{v,j})_{v\in I}\bigr)
             =\sum_{v\in I}\beta_vw_{v,j}.
\]
\end{definition}

\begin{corollary}[Recovering every Boolean witness bit]\label{cor:pack:column-bits}
For all $v\in I$ and $j\in J$,
\[
 \boxed{d_v(\widehat w_j)=w_{v,j}.}
\]
\end{corollary}
\begin{proof}
Apply Lemma~\ref{lem:pack:unpack-pack} to the
column vector $(w_{v,j})_{v\in I}$ and compare
its $v$-th coordinate.
\end{proof}

\begin{remark}[What packing does not prove]
Every field element $x\in L$ has a unique vector of
\emph{binary} coordinates in this fixed basis. Thus
the equality $x=\operatorname{Pack}(\operatorname{Unpack}(x))$
does not require a separate Booleanity constraint.
However it does not bind an independently committed
array of bits to $x$: if both representations are
provided independently, their equality is an
additional relation that must be checked.
\end{remark}

\begin{definition}[Multilinear evaluation weights]\label{def:pack:weights}
Choose two public evaluation points
$r_{\mathrm{lo}}\in L^\ell$ and
$r_{\mathrm{hi}}\in L^h$, and define
\[
 q_v:=\operatorname{eq}_\ell(v,r_{\mathrm{lo}}),
 \qquad
 h_j:=\operatorname{eq}_h(j,r_{\mathrm{hi}}).
\]
For $h=0$ the unique weight is $h_{()}=1$.
Define the true evaluation of a Boolean table by
\[
 y_w:=\sum_{v\in I}\sum_{j\in J}
             q_vh_jw_{v,j}\in L.
\]
\end{definition}

\begin{lemma}[Two-axis equality-weight factorization]\label{lem:pack:weight-factor}
For all $(v,j)\in I\times J$,
\[
 \operatorname{eq}_{\ell+h}
 ((v,j),(r_{\mathrm{lo}},r_{\mathrm{hi}}))
 =q_vh_j.
\]
\end{lemma}
\begin{proof}
The equality polynomial is the product of one
affine factor per coordinate. Its first $\ell$
factors involve only $(v,r_{\mathrm{lo}})$ and
multiply to $q_v$; its final $h$ factors involve
only $(j,r_{\mathrm{hi}})$ and multiply to $h_j$.
Splitting the finite product gives the identity.
\end{proof}

\begin{corollary}[Packed witness is the original Boolean multilinear table]\label{cor:pack:original-evaluation}
Let $\widetilde w:L^{\ell+h}\to L$ be the
multilinear extension of the field-embedded
Boolean table $w$. Then
\[
 \boxed{\widetilde w(r_{\mathrm{lo}},r_{\mathrm{hi}})
 =y_w.}
\]
\end{corollary}
\begin{proof}
Expand the multilinear extension as
$\sum_{(v,j)\in I\times J}
  w_{v,j}\operatorname{eq}_{\ell+h}
      ((v,j),(r_{\mathrm{lo}},r_{\mathrm{hi}}))$.
Substitute the factorization of
Lemma~\ref{lem:pack:weight-factor}.
The resulting sum is the definition of $y_w$.
\end{proof}

\begin{lemma}[Partition of evaluation weights]\label{lem:pack:weight-sum}
$\sum_{v\in I}q_v=1$ and $\sum_{j\in J}h_j=1$.
\end{lemma}
\begin{proof}
For each coordinate, the equality weights
for Boolean values $0$ and $1$ sum to
$(1-t)+t=1$. Distributing the finite sum
over all bit strings as a product of these
two-term sums gives $1$ for each axis.
When $h=0$, the second sum has one term $1$.
\end{proof}
